%
%! regex-fix=true
%! dir=rataflechera

\songtitle{Draw the Line}
\tag*{2026}\tag{rock-n-roll}
\begin{notes}
Lyrics by Suno based on a chapter written by Claude summarizing several original essays by Carlos Thompson -- a case against false moral equivalence: two wrongs can both be wrong without being the same wrong, and power imbalance changes the weight of a label. Retrieved from rataflechera's Suno page (V6, 29 September 2026).
\end{notes}
\begin{style}
Rock-n-roll with punchy live-kit drums, gritty overdriven electric-guitar chords, melodic bass, and restrained analog synth accents; driving straight-ahead groove, lean verses opening into broad forceful choruses, with a brief instrumental guitar-synth hook between sections and a firm band stop. Clear, focused lead vocal with crisp chorus layering; clean, polished 2000s production and energetic mix.
\end{style}
\begin{lyrics}
[Verse 1]
A scratch across the fender, a life beneath the wheel
Both can cross the line, but they don't carry equal steel
Call the harm by its name, don't flatten what it means
A mark against the paint is not a body in the street

[Pre-Chorus]
The measure isn't neutral
When the weight is on one side

[Chorus]
Draw the line, keep it clear
Name the force that's standing here
Two wrongs can both be wrong
Without pretending they're the same
Draw the line, hold it fast
Ask whose power makes the last word last
No clean hands, no easy frame
Still, unequal isn't equal all the same

[Verse 2]
A station set to burning, a bull brought to its knees
A citizen struck down by the hand sworn to keep peace
A mosquito crushed at dusk; the word can cover all
But a label isn't judgment, and a label isn't law

[Pre-Chorus]
The office changes everything
The duty makes a claim

[Chorus]
Draw the line, keep it clear
Name the force that's standing here
Two wrongs can both be wrong
Without pretending they're the same
Draw the line, hold it fast
Ask whose power makes the last word last
No clean hands, no easy frame
Still, unequal isn't equal all the same

[Bridge]
Not every act of order earns the name of right
Not every act of anger turns to justice overnight
Don't hand the state a blank check
Don't bless each broken pane
Keep the facts in proper focus
Let the hardest questions stay

[Final Chorus]
Draw the line, keep it clear
Name the force that's standing here
Two wrongs can both be wrong
Without pretending they're the same
Draw the line, hold it fast
Ask whose power makes the last word last
No clean hands, no easy frame
Still, unequal isn't equal all the same

[Outro]
Name it true
Weigh it through
Unequal isn't equal all the same
\end{lyrics}

%! dir=chlewey

\songtitle{The Same Colombia}
\tag*{2026}\tag{folk-rock}\tag{original-source}\tag{thomas-entourage}\tag{erika-czajka}\tag{victoria-ospina}\tag{natalia-kyou-park}\tag{claire-jones}\tag{thomas-garcia}
\begin{notes}
Lyrics by ChatGPT based on a short story written by Carlos Thompson -- in the voice of Erika Czajka, travelling through Colombia with Vicky for the first time: Punta Espada, Natalia's news about Claire and Thomas, and the thing Erika finds herself imagining but never saying out loud. Retrieved from Carlos Th's Suno page (V6-MINI, 29 September 2026).
\end{notes}
\begin{style}
Folk rock with Swedish visa inflections; female mezzo vocals, opening in intimate spoken narration before turning to a melodic chorus; strummed acoustic guitar, bowed bass, accordion melody, muted electric guitar, synth-pad swells and subtle analog synth arpeggios; 92 BPM with a swung eighth-note feel; plate reverb, tape saturation and stereo delay throws.
\end{style}
\begin{lyrics}
[Intro - Spoken]
Vicky and I in Colombia- what could go wrong?

[Verse 1]
Punta Espada is now feeling real,
A nice resort along the shore,
Separate rooms because of optics,
Though the staff seemed not to care much more.
Half work, half vacation,
For Vicky and for me,
Walking through a place I know
As if I'd never seen.

[Pre-Chorus 1]
And there was something in the air,
A little thrill, a little dare,
The kind of thought you don't confess
Because you know where it could lead.
Nothing had happened, nothing yet,
But sometimes that's when danger starts--
The anticipation of transgression
Before anybody crosses the line.

[Chorus]
Our first time travelling together,
Through a country we both somehow know,
Familiar streets and unfamiliar faces,
Places known, but seen anew.
The same Colombia beneath our feet,
The same old roads, the same old heat,
Yet everything was strangely new
Because I was there with you.

[Verse 2]
I didn't come to the inauguration.
Natalia said it was discreet,
Beautiful and emotive,
A sense of duty made complete,
Against the backdrop of tensions
And the tragedy next door.
Then she told me something else
I wasn't prepared for.

[spoken:] Claire.

I could feel Claire and Thomas
Were into something more,
But I always thought it was platonic,
Something I'd seen before.
Vicky never suspected,
Then mildly shocked, she said
That suddenly some little things
Were making sense inside her head.

[Pre-Chorus 2]
And on the last flight they weren't
Indeed pretending anymore.
I kept thinking it was eerie,
Like a door left open in a wall.
But Vicky pointed out to me
That I always react that way
When Thomas gets that physically close
To anybody anyway.

[Chorus]
Our first time travelling together,
Through a country we both somehow know,
Familiar streets and unfamiliar faces,
Places known, but seen anew.
The same Colombia beneath our feet,
The same old roads, the same old heat,
Yet everything was strangely new
Because I was there with you.

[Bridge]
What troubled me the most
Wasn't Claire, or what they did.
It was that Thomas never told me,
And I wondered why he hid it.
Maybe secrets feel much larger
When you're travelling together,
When the distance makes you notice
Every thread between each other.

[Final Chorus]
Then Bogotá, a small departure
From the usual circuit we knew,
Visiting Victoria's aunts and uncles,
The open-minded and the not-so-much too.
Always introduced as Vicky's friend,
Something I had never cared about,
Until I found myself imagining
What it would feel like to say it out loud.

[Outro]
And Vicky was introduced as my friend
To Francine and Gabriel,
In a way that made the texture
Of our friendship clear as well.

We walked the evening streets
Like ordinary tourists,
Through a country equally familiar,
And equally strange.

[Spoken]
What could go wrong?

Nothing.
\end{lyrics}

%! dir=wyomee

\songtitle{Under Aachen Sky}
\tag*{2026}\tag{medieval-epic}\tag{original-source}\tag{african-empire}\tag{monukae}\tag{mombuto}\tag{gladiel}
\begin{notes}
Adaptation by Claude of a dialog written by Carlos Thompson -- Mombuto and Gladiel's wedding under Charlemagne's Aachen, reconciling a history of conquest, exile, and clashing faiths into one shared vow. Retrieved from Wyomee Dank's Suno page (V6-MINI, 29 September 2026).
\end{notes}
\begin{style}
Medieval epic ballad
\end{style}
\begin{lyrics}
[Verse 1: Mombuto - male tenor]
I could not fathom why my mother prayed
To sun and peak in Monukae's fierce heat,
Till in the Andes, where her fathers strayed,
The heights made plain what once was incomplete.

[Verse 2: Gladiel - female soprano]
I could not pardon God, who let the horde
From desert wastes lay waste my native town,
Who slew my kin, and set me at their board,
And made my brother bring my father down.

[Verse 3: Mombuto]
The jungle gods my grandsire brought from far,
The monks' dear Christ, the scholars' Prophet true,
My father laughed at all, and by that star
I learned to hold all faiths in equal view.

[Verse 4: Gladiel]
Thy father taught me every god to prize,
And helped my heart make peace with mine at last;
He mocked, yet honoured; and within thine eyes
I know that same respect is holden fast.

[Chorus: both]
For gods may differ, yet the heart is one;
Through fire and exile love at last is won.

[Verse 5: Mombuto]
So proud thou wert, so lofty, though brought low,
I could not choose but look on thee and yearn;
Thy scorn of me did make thy beauty grow
The brighter in mine eyes, and made me burn.

[Verse 6: Gladiel]
I saw thee come to Cork: a comely sight,
I thought, and let that thought no further stray;
So long I could not honour thee aright,
For all the pain that followed on that day.

[Verse 7]
(Mombuto:) Thou know'st I am a man of words, not steel;
(Gladiel:) I know it now; but then I could but feel
That for thy sake my world was made to kneel.

[Intermezzo - instrumental verse]

[Verse 8]
(Gladiel:) Then from beyond the sea thou cam'st once more;
Not Marco, but with Numak at thy side.
(Mombuto:) Well was young Marco when I left that shore;
He serves the Inca for my sire with pride.

[Verse 9]
(Gladiel:) Thy father told me true, yet blame I chose.
(Mombuto:) Yet now thine eyes no longer hate disclose.

[Verse 10]
(Gladiel:) No more; I understand. Take thou my frame,
Take thou my God; let love undo the shame.
(Mombuto:) I take thee, Gladiel, and take thy God;
Let Aachen's bishop seal it, by his rod.
(Gladiel:) I take thee, Mombuto, and thee alone;
Let bishop and the Frankish emperor own
Our love, and make thy heart and mine as one.

[Final Chorus: both]
For gods may differ, yet the heart is one;
Through fire and exile love at last is won.

[Final Chorus: both]
For gods may differ, yet the heart is one;
Through fire and exile love at last is won.
\end{lyrics}

%! dir=dvigitt

\songtitle{When I Say Yes}
\tag*{2026}\tag{r-and-b}\tag{original-lyrics}\tag{thomas-entourage}\tag{natalia-kyou-park}\tag{thomas-garcia}
\begin{notes}
Written by Carlos Thompson, in the voice of Natalia Kyou Park from \emph{Thomas' Entourage} -- her goodbye to their arrangement as she commits to Michael, one last time before she says yes. Retrieved from Dvigitt's Suno page (V6-MINI, 29 September 2026).
\end{notes}
\begin{style}
R\&B pop with candy-pop sheen, intimate close-mic production, plate reverb and stereo delay throws; breathy female vocals, spoken verses, stacked harmonies and a call-and-response chorus; syncopated claps, elastic bassline, glossy synth plucks and soft electric piano; midtempo groove with sidechain pumping.
\end{style}
\begin{lyrics}
[Verse 1]
Hello, Thomas. I don't know what you'd think.
You seem the kind of person who could understand.
Our little arrangements must come to an end.
Michael and I are sealing our arrangement in ink.

[Pre-Chorus]
My old ways now seem wrong.
Michael had Los Feliz;
I had the rest of the world for my own.

[Chorus]
But when I say yes,
it will feel wrong to fool around.
The world won't be mine any longer;
my world will be his now,
and I know you'll understand.

[Verse 2]
Hello, Thomas. It is time to say goodbye.
I am still at your professional service,
but the thing within us shall not resurface
one last time, and then just let it die.

[Pre-Chorus]
One more time to make it wrong,
then make it right for Michael.
Say yes, and then I'll own what I'll own.

[Chorus]
Then, when I say yes,
it will feel wrong to fool around.
The world won't be mine any longer;
my world will be his now,
and I know you'll understand.

[Bridge]
You're the one I'll choose
to end this cycle of my life,
'cause I know you'll understand.
You're the right one to make it wrong
one last time before I make it right,
one last time before I decide
if I want you and Gabi
at my wedding, at my side.

[Final Chorus]
'Cause when I say yes,
it will feel wrong to fool around.
The world won't be mine any longer;
my world will be his now,
and I want you to share this with me
'cause I know you'll understand.
\end{lyrics}

%! dir=gabisson

\songtitle{Worthy}
\tag*{2026}\tag{sophisti-pop}\tag{original-lyrics}\tag{thomas-entourage}\tag{claire-jones}\tag{thomas-garcia}
\begin{notes}
Written by Carlos Thompson, in the voice of Claire Jones from \emph{Thomas' Entourage} -- reassuring Thomas that his worth isn't his bank account or his self-doubt, but the people, entourage included, who stay anyway. Retrieved from Gabriel's Suno page (V6-MINI, 29 September 2026).
\end{notes}
\begin{style}
Sophisti-pop, Welsh art pop with Rhodes electric piano, chorus guitar arpeggios, fretless bass lines, muted synth brass, brushed live-kit groove and handclap backbeat; steady 92 BPM; plate reverb, tape saturation and widened stereo; intimate female vocals with clear, warm phrasing and a gently lifting chorus.
\end{style}
\begin{lyrics}
[Verse 1]
You must understand that you're worthy on your own, my love.
You know you can achieve greatness, and you've done so,
but somehow you still believe your human connections are at fault,
despite me at your side, despite the people who are still on.

[Pre-Chorus]
Many wish for the texture of your relationship with your ex,
the unconditional partnership of Gabriele in your life,
the loyalty of Pilar, Susana, Fernando, Sandra, and myself,
and the love you know I still profess to you, my dear.

[Chorus]
You're worthy, my love.
You move things, you achieve greatness.
You've made life so good for so many.
You're worthy, my love.
Even when you believe your life's a mess,
you've got an entourage of caring.

You're worthy, my dear.

[Verse 2]
Your worth is not your bank account, and you know it, my love.
People who remain at your side are not interested in your wallet.
Those who are soon understand they won't survive with you.
And you should not be ashamed that they are just gone.

[Pre-Chorus]
Your parents, your son, your daughter: that's real love.
Your friendship with Erika, and Fiona, and Xumi,
whatever you and Diana are playing, all that's very real.
And that's a reason they and I are still around you, my dear.

[Chorus]
You're worthy, my love.
You move things, you achieve greatness.
You've made life so good for so many.
You're worthy, my love.
Even when you believe your life's a mess,
you've got an entourage of caring.

You're worthy, my dear.

[Bridge]
And you make life worthy for so many--
that's one of the things I'm drawn to in you.
You have a clear vision, humble but sharp;
you understand limits like few do with your reach.
You know you're more than what your mind
keeps lying you're not, my dear.

[Final Chorus]
You're worthy, my love.
You move things, you achieve greatness.
You've made life so good for so many.
You're worthy, my love.
Even when you believe your life's a mess,
you've got an entourage of caring.

You're worthy, my dear.

Worthy.
\end{lyrics}

%! dir=chlewey

\songtitle{Leading the Jump}
\tag*{2026}\tag{blues-rock}\tag{original-lyrics}
\begin{notes}
Written by Carlos Thompson -- 1995 to 1997, joining a systems-engineering R\&D group and juggling teaching-assistant work in math and programming, pioneering alongside the internet's own early opening-up, and a string of crushes that all led nowhere until the one that finally did. Retrieved from Carlos Th's Suno page (V6-MINI, 30 September 2026).
\end{notes}
\begin{style}
Blues rock, 1990s alt rock; tape saturation, plate reverb, stereo delay throws; electric guitar riffs and slide guitar licks over an overdriven bassline, live drums with tom-heavy fills and organ stabs; baritone lead with layered chorus harmonies in call-and-response; 132 BPM, swung backbeat, dynamic breakdowns.
\end{style}
\begin{lyrics}
[Verse 1]
I applied to a group
in research and development
on systems engineering
as an electronics student.

I knew some programming--
C++ and C.
I was qualified,
and they adopted me.

[Pre-Chorus]
R&D in computing,
teaching assistant in math,
TA in programming--
I was a busy mind.

[Chorus]
Nineteen ninety-five,
Internet opening up.
I felt like pioneering,
leading the jump.

[Verse 2]
A new R&D approach--
I followed through.
Our leader insisted
we were social folk too.

Engineering classes,
teaching assistance,
computing R&D,
and less social distance.

[Pre-Chorus]
And right next to our office,
another student group--
industrial eng. babes,
a new partying troop.

[Chorus]
Nineteen ninety-six,
Internet opening up.
They called me an expert,
leading the jump.

[Verse 3]
For me the word "crush"
has always been plural,
an easy position
in a closed journal.

Crushes added up
with new industrial friends.
Add them to the list
of my many dead ends.

[Bridge]
And that was the state
of my affairs:
feeling so close,
my heart in crosshairs.

[Chorus]
Nineteen ninety-seven,
Internet opening up.
That was my edge,
leading the jump.

[Outro]
And my life changed
from parties and longing
to someone who brought me
to a new way of belonging.
\end{lyrics}

%! dir=rataflechera

\songtitle{Mueca Pantera}
\tag*{2026}\tag{ambient-dubstep}\tag{original-lyrics}\tag{spanish-lyrics}
\begin{notes}
Written by Carlos Thompson, sung in his own recorded voice -- a short a cappella chant for a black cat stretching and yawning in a hammock by the window. Retrieved from rataflechera's Suno page (V6, 30 September 2026).
\end{notes}
\begin{style}
Ambient dubstep, the user's own recorded baritone voice in a cappella chant over a subbass drone; hyper-compressed production; slow, spacious pulse with heavy dubstep subbass drops; opening death-metal riff, then synthesized new age melody.
\end{style}
\begin{lyrics}
[Intro - Death metal riff in D-minor]

[Verse - baritone vocals, a capella over a subbass drone]
Ese bostezo,
mi mueca pantera,
en tu hamaca
frente a la ventana.

Esa estirada,
panterita negra,
imutado en tu sueño
¡qué vida esa!

[Outro - new age, ambient dubstep]
\end{lyrics}

%! dir=gabisson

\songtitle{Två år på Söder}
\tag*{2026}\tag{vispop}\tag{swedish-lyrics}\tag{two-years-in-sodermalm}
\begin{notes}
Swedish-language adaptation by Claude of \emph{Two Years in Södermalm} (page \pageref{sec:two-years-in-sodermalm}), the same two years of noticing a stranger on the commute past Medborgarplatsen to Gullmarsplan's Vuxengymnasium. Retrieved from Gabriel's Suno page (V6-MINI, 30 September 2026).
\end{notes}
\begin{style}
Vispop, keyboards and electric guitars underpin a steady early-1990s groove, with fiddle filling the melodic spaces; early-90s digital reverb; male baritone lead, warm and direct.
\end{style}
\begin{lyrics}
[Vers 1]
Jag såg dig på avstånd i snart två år,
kring Medborgarplatsen, där tunnelbanan går,
på väg till vuxengymnasiet vid Gullmarsplan,
ibland på andra håll i stan.

[Vers 2]
Jag fick aldrig veta vem du var, aldrig ditt namn,
bara din stil, samma upproriska punk varje säsong,
ändå bara en fast punkt i stadens sammanhang.

[Refräng]
Två år på Söder, jag gick och jag gick,
en fast punkt i människornas landskap:
en rödhårig punkare, en rödtott med blick,
som jag nästan blev kär i.

[Vers 3]
Du höll mig förtrollad i snart två år,
två år såg jag dig i kvarteren där vi går,
i ett kvarter för stort för att du skulle se mig,
länets skeva mitt, där ingen ser varandra.

[Refräng]
Två år på Söder, jag gick och jag gick,
en fast punkt i människornas landskap:
en rödhårig punkare, en rödtott med blick,
som jag nästan blev kär i.

[Brygga]
Samma jacka, samma strumpor, året runt,
samma störande punkestetik,
en mask över en vacker kvinna därinne
som jag inte kunde låta bli att längta efter.

[Sista refräng]
Två år på Söder, jag gick och jag gick,
en fast punkt i människornas landskap:
en rödhårig punkare,
en vacker kvinna under,
ett långt längtande, en lång dröm,
en flicka som jag nästan blev kär i.
\end{lyrics}

%! dir=wyomee

\songtitle{The Fixed Point}
\tag*{2026}\tag{soft-rock}\tag{original-source}\tag{thomas-entourage}\tag{thomas-garcia}\tag{pilar-venegas}\tag{diana-naji}\tag{aurora-swansen}\tag{susana-bianchi}\tag{sylvia-lai}
\begin{notes}
Lyrics by Claude based on an original text by Carlos Thompson, corrected and expanded by ChatGPT -- Thomas rebuilding after the Pereira landslide while Pilar stays the fixed point everyone else moves around, and Diana unexpectedly finding she loves a seat on the family board. Retrieved from Wyomee Dank's Suno page (V6-MINI, 30 September 2026).
\end{notes}
\begin{style}
Soft rock ballad, intimate male vocal, unhurried 80 BPM sway; piano and acoustic guitar lead, delayed electric-guitar swells, warm understated bass, brushed drums, light congas and güiro, subtle dembow entering only in the final chorus; warm layered backing vocals in choruses, close and spacious mix.
\end{style}
\begin{lyrics}
[Verse 1] (piano or nylon guitar, close vocal)
The penthouse lights are dark in Bogotá,
He's in the cabin, listening to the sea.
Since the ground came down in Pereira,
He's been out where the work needs to be.
Trucks on the road, cranes by the station,
Steel going up where the walls came down.
Months of rebuilding, one slow foundation,
Quiet man in a half-built town.

[Pre-Chorus] (soft electric guitar enters)
The dust has settled, the cranes are slow,
Somebody's got it, he doesn't have to know.

[Chorus]
Everyone's moving, everyone's turning,
Some are leaving, some returning,
And Pilar stays, the way she's always stayed,
The fixed point where it all gets arranged.

[Verse 2: Diana]
Ilia's health has shaken the house,
The family firm needs a steady hand.
Diana, who was never home,
Now joins the board from wherever she lands.
Still on the road, but she calls in now,
Numbers at dawn, the board on screen.
And the thing that no one saw coming:
She loves it, in a way no one's seen.

[Verse 3: the rest of them]
Aurora's counting down to October,
Susana's out in the desert sand,
Sylvia's got a name on her own door now,
And none of it was ever planned.

[Bridge] (strip back to guitar and voice, then slow build)
Some lives are updates, some lives are weather,
Some lives are the ground you stand upon.
She doesn't call, she doesn't gather,
She's just there when the dust is gone.
He looks at the hills, and the hills look back,
And nothing is waiting for him to act.

[Final Chorus] (full band, layered harmonies, dembow pulse enters quietly under the second line)
Everyone's moving, everyone's turning,
Some are leaving, some returning,
And he stays, the way she's always stayed,
And for once, it moves, unarranged.

[Outro] (dembow fades, guitar and voice only)
For once, unarranged...
(Just the sea.)
\end{lyrics}

\songtitle{Public Knowledge}
\tag*{2026}\tag{soft-rock}\tag{original-source}\tag{thomas-entourage}\tag{thomas-garcia}
\begin{notes}
Lyrics by Claude based on an original text by Carlos Thompson, corrected and expanded by ChatGPT -- the open secret of Claire and Thomas, alongside Mariain's non-exclusive painter, Xumi and Pere Joan's undefined thing, and Natalia and Michael's wedding date: all the arrangements everyone knows and nobody names. Retrieved from Wyomee Dank's Suno page (V6-MINI, 30 September 2026).
\end{notes}
\begin{style}
Soft rock with dry, close-miked vocals; relaxed 92--98 BPM dembow sway; clean palm-muted electric guitar, warm bass, soft synth pads, restrained percussion, and a lighter half-time chorus pulse; conversational half-sung verses rise into open-chord choruses with layered harmonies, then a percussion-only bridge builds to a full-band final chorus and guitar-led outro.
\end{style}
\begin{lyrics}
[Verse 1]
Claire and Thomas said it out loud,
Only to the ones who need to know.
Out in public, easy smiles,
Gabi's name still on the show.
Vicky's in London, Erika's in Utrecht,
Two front doors and one new word:
They quit calling it "just friends,"
Though nobody's said it, nobody's heard.

[Pre-Chorus]
Everybody's got a story,
Nobody's got the name.
Say it low, say it slow,
Keep it in the family.

[Chorus]
We don't have a word for it yet,
Everyone knows, but nobody's said.
Those who know won't tell a soul,
It's public knowledge, so keep it low.

[Verse 2]
Mariain's got a painter from Arkansas,
"It's not serious," she swears.
"Non-exclusive," says it twice,
Which, coming from her, means she's the one who decides.
Xumi and Pere Joan, nothing official,
Neither one will draw the line.
When nobody wants to name it,
That's how you know it's not just fine.

[Pre-Chorus]
Everybody's got a story,
Nobody's got the name.
Say it low, say it slow,
Keep it in the family.

[Chorus]
We don't have a word for it yet,
Everyone knows, but nobody's said.
Those who know won't tell a soul,
It's public knowledge, so keep it low.

[Bridge] (percussion only, then build)
Natalia and Michael set the date,
March, and the invitations are out.
Erika plus one, Thomas plus one,
Who's a friend and who's a secret,
Who was told and who's still in doubt?
Raise a glass, pass the bread,
Remember what the room can know.
Everyone's smiling, everyone's counting,
Nobody's saying so.

[Final Chorus] (full band, layered harmonies)
We had the word all along,
Soft as a secret, steady and strong.
Now you know, and it's no surprise,
It's public knowledge in everyone's eyes.

[Outro] (guitar and hummed hook)
We don't have a word...
(We never needed one.)
\end{lyrics}

\songtitle{Somewhere Between Flights}
\tag*{2026}\tag{pop-rock}\tag{original-source}\tag{thomas-entourage}\tag{sylvia-lai}\tag{aurora-swansen}\tag{zahra-moradi}\tag{jan-alan}\tag{mihai-rusu}
\begin{notes}
Lyrics by Claude based on an original text by Carlos Thompson, corrected and expanded by ChatGPT -- Sylvia settled into her own name, Gabriele splitting her life between Potsdam, Queens and Poblenou, Fiona and Susana each finding their own geography, and Aurora and Zahra moving back through an orbit that keeps changing without them. Retrieved from Wyomee Dank's Suno page (V6-MINI, 30 September 2026).
\end{notes}
\begin{style}
Melodic pop rock; tape saturation, plate reverb, and sidechain-pumped swells; electric guitar arpeggios, clean rhythm guitar, synth bass ostinato, piano chords, analog pads, tom-heavy percussion, muted congas, and stereo hand percussion, with a bridge percussion drop and pre-chorus lift; light dembow groove at 100 BPM; female lead with layered harmonies.
\end{style}
\begin{lyrics}
[Verse 1: Sylvia]
She hung the diploma, made no fuss,
Just took her seat like it was always hers.
Used to orbit, now she holds her own,
Signing her name to the things she's earned.
Inside the system, learning the rules,
Keeping some, bending a few.
Same Sylvia, same laugh at the door,
Only now she's the one who's grown.

[Pre-Chorus]
Gate changed, flight delayed,
Nobody's lost, they're just on their way.

[Chorus]
We're somewhere between flights,
Between the lease and the studio lights,
Home is a city you call by its name,
And it never quite stays the same.
Ask me where I am tomorrow,
I'll tell you where I'll be,
Somewhere between flights,
And happy as can be.

[Verse 2: Gabriele]
A studio in Potsdam, signed and sealed,
A share of the lease still holding in Queens.
But Poblenou's where the paint dries now,
Where the local scene finally made her welcome.
Twice a month she goes back east,
Enough to keep the old room warm.
She doesn't live there, she doesn't leave,
She just lets it breathe, and then she's gone.

[Pre-Chorus]
Gate changed, flight delayed,
Nobody's lost, they're just on their way.

[Chorus]
We're somewhere between flights,
Between the lease and the studio lights,
Home is a city you call by its name,
And it never quite stays the same.
Ask me where I am tomorrow,
I'll tell you where I'll be,
Somewhere between flights,
And happy as can be.

[Verse 3: Fiona and Susana]
Fiona's in San Jose, it's on her usual round,
But the Bay was never just a stop.
It's where she learned to build things right,
Where her mind took shape, and the code took root.
Susana found a desert she loves,
The Guajira, boots on, wind in her hair.
Half in Punta Espada, half in Bogotá,
Of course she hikes the desert, between the flights.

[Bridge: Aurora, then Zahra] (drums drop to percussion, build slowly)
Aurora's been away two months,
Counting the days like an old calendar.
Mid-October, Bogotá, home again,
The city noticed she was gone.
Then December, Punta Espada,
Zahra's flying in with Michai and Jan-Alan.
Back in the orbit she left for a while,
She'll see how it's changed, and she'll understand.

[Final Chorus] (full band, layered harmonies)
We're somewhere between flights,
But the ones who stay have the lights,
Someone's always at the gate,
Someone's always arriving late.
Ask me where I'll be in December,
I'll tell you where I'll be,
Somewhere between flights,
But the road leads back to me.

[Outro] (guitar and soft percussion, voice fading)
Gate changed, flight delayed,
Nobody's lost, they're just on their way.
\end{lyrics}

%! dir=dvigitt

\songtitle{I Chose This Life}
\tag*{2025}\tag{adult-contemporary}\tag{original-source}\tag{thomas-entourage}\tag{xumi}
\begin{notes}
Lyrics by Suno based on the closing chapter of Xumi's memoir, co-written by Carlos Thompson and ChatGPT (2025) -- Xumi's own account of the Barcelona life she built on her own terms, friends Pere Joan and Juan Carlos included, choosing it rather than waiting for anything else. Retrieved from Dvigitt's Suno page (V6-MINI, 30 September 2026).
\end{notes}
\begin{style}
Adult contemporary with acoustic guitar riffs shaped by pasodoble: a steady, measured march pulse softened into a warm mid-tempo pop groove, bright Spanish acoustic strums answering the vocal, restrained bass and brushed percussion, subtle piano support, verses kept close and conversational, chorus opening into broader harmony and fuller rhythm, then a reflective guitar-led outro. Female mezzo-soprano, clear and grounded with natural phrasing; polished, spacious mix with organic dynamics.
\end{style}
\begin{lyrics}
[Verse 1]
Pere Joan came through friends of friends
In twenty-twenty-three, Barcelona wide
A handball player, steady in a contest
A friend who never asked me to revise
We watched the matches, weighed the tactics
Handball or futsal, which had changed the game?
He knew the lines I drew between us
And never tried to move them all the same

[Chorus]
This is not a waiting room
Not a door I'm counting on
I made a life that fits my hands
And I wake up inside it strong
Good friends, a table, work that matters
A match at hours the clocks refuse
I'm not holding out for something else
I chose this life, and I choose

[Verse 2]
Juan Carlos came the following year
From Gràcia, where his restaurant drew me in
He made a welcome feel like breathing
No need to watch the way I settled in
He cooked; we talked about the football
His Venezuelan view took patient schooling
Warm and direct, with room for laughter
No hidden claim, no borrowed rule

[Chorus]
This is not a waiting room
Not a door I'm counting on
I made a life that fits my hands
And I wake up inside it strong
Good friends, a table, work that matters
A match at hours the clocks refuse
I'm not holding out for something else
I chose this life, and I choose

[Bridge]
Eixample mornings, uncluttered shelves
Economics books I read for real
A good sound system, work worth doing
The quiet pride of knowing what I feel
Recife's score comes through the distance
Joy and ache on a European clock
And across the cities, friends are there
Because we chose each other well

[Final Chorus]
This is not a waiting room
There's no missing life to come
I built this one from what was true
And it's already more than enough
Good friends, a table, work that matters
A match at hours the clocks refuse
I'm not holding out for something else
I chose this life, and I choose

[Outro]
I chose this life
And I choose
\end{lyrics}

%! dir=gabisson

\songtitle{A Ground I Choose}
\tag*{2026}\tag{folk-pop}\tag{gabriele-wiater}\tag{thomas-entourage}\tag{thomas-garcia}
\begin{notes}
Lyrics by Suno based on a ChatGPT conversation on why Gabi would not choose to marry Thomas for an EB-5 visa -- her “ground” is a life that does not depend on his name or his papers. See \emph{Leave the Door Unlatched}. Retrieved from Gabriel’s Suno page (V6-MINI, 2 October 2026).
\end{notes}
\begin{style}
Folk pop with industrial undertones, fingerpicked acoustic guitar over muted metal percussion and low analog bass, measured mechanical pulse, restrained verses opening into a firm melodic chorus, dry tactile mix with subtle room depth, intimate close-miked female alto.
\end{style}
\begin{lyrics}
[Verse 1]
The lawyer laid the papers straight
Beside the coffee gone cold
Said, “There’s a road that could hold you here,”
And named the risks in careful prose

You watched the window, counted costs,
The years you’d spent staying light
A key that fits another door
Can still be heavy in your hand

[Pre-Chorus]
You know the promise, know the price
Know what the fine print can’t define

[Chorus]
I won’t build my life on “if you stay”
Or place my name where yours can sway
A safer road can close around me
A steady hand can still let go
I want a ground that answers only
To the steps I choose to own

[Verse 2]
You know his maps cross many borders
His days don’t settle in one place
He could learn the forms and show up
But you can read the strain behind his face

You’ve kept your love in honest words
No office stamp, no borrowed frame
A vow can make a home feel smaller
When the world records it under one name

[Pre-Chorus]
You won’t call fear a lack of trust
You’ve just learned what must be yours

[Chorus]
I won’t build my life on “if you stay”
Or place my name where yours can sway
A safer road can close around me
A steady hand can still let go
I want a ground that answers only
To the steps I choose to own

[Bridge]
Barcelona, a room with my own key
Morning bread, a desk, the sea
No promise borrowed for a season
No future balanced on his being

[Final Chorus]
I won’t build my life on “if you stay”
I know the cost; I know the way
The strongest door is not the one
That needs another hand to close
I want a ground that answers only
To the life I choose to hold

[Outro]
The papers rest beside the cup
She takes her coat and walks out sure
\end{lyrics}

\songtitle{Leave the Door Unlatched}
\tag*{2026}\tag{folk-pop}\tag{gabriele-wiater}\tag{thomas-entourage}\tag{thomas-garcia}
\begin{notes}
Lyrics by Suno based on a ChatGPT conversation on why Gabi would not choose to marry Thomas for an EB-5 visa, sung to Thomas rather than about him. See \emph{A Ground I Choose}. Retrieved from Gabriel’s Suno page (V6-MINI, 2 October 2026).
\end{notes}
\begin{style}
Folk pop, restrained and intimate, with a measured, lightly pulsing groove; close-miked female lead, conversational in spare verses and opening into layered chorus harmonies; fingerpicked acoustic guitar leads muted piano, found-metal percussion and low mechanical thuds, with warm bass and clipped electric textures building gradually; bridge falls back to voice and guitar, final chorus fuller but controlled; dry upfront vocal, tactile room detail, deliberate dynamics, polished organic mix.
\end{style}
\begin{lyrics}
[Verse 1]
You laid the papers on the table
Beside the keys and coffee rings
A clean route through a crowded doorway
A promise folded into things

You said the file could hold our story
The dates, the trips, the years we chose
I heard the care inside your answer
And all the doors it would close

[Pre-Chorus]
I know it could work
I know what it asks
A home made official
A life built to last

[Chorus]
I won’t make your name my shelter
I won’t make my ground depend
On the version of us that never falters
On a road that bends around one man
You can love me without holding
Every piece of where I stand
I can leave the door unlatched between us
And still know exactly who I am

[Verse 2]
You move with maps in several pockets
Your days can cross an ocean wide
I’ve learned the shape of your devotion
And where your restless seasons hide

It isn’t doubt that makes me careful
Or fear of what we might become
I want a life that keeps its balance
Even when the plans come undone

[Pre-Chorus]
I know it could work
I know what it asks
A home made official
A life built to last

[Chorus]
I won’t make your name my shelter
I won’t make my ground depend
On the version of us that never falters
On a road that bends around one man
You can love me without holding
Every piece of where I stand
I can leave the door unlatched between us
And still know exactly who I am

[Bridge]
Let the forms stay in their folder
Let the future keep its room
I would rather choose the distance
Than make a promise out of proof

[Final Chorus]
I won’t make your name my shelter
I won’t make my ground depend
On the version of us that never falters
On a road that bends around one man
You can love me without holding
Every piece of where I stand
I can leave the door unlatched between us
And still know exactly who I am

[Outro]
I know what you offered
I know what it meant
I’m keeping my own way open
And I’m not calling that an end
\end{lyrics}

%! dir=rataflechera

\songtitle{A Narrow Way}
\tag*{2026}\tag{chamber-pop-rock}\tag{ricochets}\tag{fabian-muller}
\begin{notes}
Lyrics by Suno based on a Claude explanatory summary of the conversation between Müller and Sobal in \emph{Ricochets}, chapter 1. Retrieved from rataflechera’s Suno page (V6, 2 October 2026).
\end{notes}
\begin{style}
Chamber pop-rock with a measured, tense walking pulse; clipped piano and muted electric guitar under close, conversational lead vocals, supported by restrained strings. Build from spare verses into a broad, controlled chorus with cello counterlines and firm live drums. Keep the mix intimate and dry at first, widening in the chorus; bridge shifts into crisp, articulate rap over a lean bass-and-percussion groove, then returns to a composed final chorus.
\end{style}
\begin{lyrics}
[Verse 1]
Two crews in the lobby, no attempt to disappear
A bald beard at his shoulder, polished shoes at either ear
No one raises a voice; the elevator doors divide
One look from Müller, and the meeting steps outside

[Pre-Chorus]
Away from the room where Durand sleeps
A quiet turn, a boundary drawn
The street receives the terms he keeps
Before the first complaint is gone

[Chorus]
Keep your voice low, let the whole block hear
Who can move the meeting, who can make it clear
A name held formal, a first name laid down
One hand on the door, one eye on the town
No promise in the pastry, no peace in the smoke
Just a narrow way through the words that we spoke

[Verse 2]
Sobal takes his time with the cigarette flame
“We haven’t got what we were promised”—plain blame
Müller points to Mendoza, says, “That’s his department”
Sobal says “Mr. Fabian Muller”—each syllable deliberate

[Pre-Chorus]
“Barys,” says Müller, softer now
Then scans the corners, roofline, glass
“You wouldn’t want to talk with me right now”
A warning shaped to let him pass

[Chorus]
Keep your voice low, let the whole block hear
Who can move the meeting, who can make it clear
A name held formal, a first name laid down
One hand on the door, one eye on the town
No promise in the pastry, no peace in the smoke
Just a narrow way through the words that we spoke

[Bridge - Rap]
Two sides arrive with their leverage displayed
The cost of a threat is already paid
So don’t spend it twice, don’t sharpen the scene
Use process as cover, keep the edges clean
Name him formally: “I know who you are”
Then answer with “Barys”—bring the distance less far
Claim you’re being watched; let the rooftops reply
Maybe it’s true, maybe useful to imply
He says, “Gunther kept his word”—now the wound is trust
Not the shipment, not the timetable, not the dust
Offers the case; Müller smiles, lets it stay
Refuses the smoke without pushing away
“Mendoza is what I can offer.” Door gently shut
A bakery, warm bread, coffee in a cup
Walk the last block without making a show
Make the ordinary where the hard feelings go

[Final Chorus]
Keep your voice low, let the whole block hear
Who can move the meeting, who can make it clear
A name held formal, a first name laid down
One hand on the door, one eye on the town
No promise in the pastry, no peace in the smoke
Just a narrow way through the words that we spoke

[Outro]
Coffee, bread, and a measured goodbye
The street keeps walking; the answer stays dry
\end{lyrics}

%! dir=wyomee

\songtitle{Every Fact in Its Own Time}
\tag*{2026}\tag{hip-hop-pop}\tag{ricochets}\tag{david-harris}\tag{marion-gebhart}\tag{john-carlson}\tag{lilith}
\begin{notes}
Lyrics by Suno based on a ChatGPT narrative profile of David Harris in \emph{Ricochets}. Retrieved from Wyomee Dank’s Suno page (V6-MINI, 2 October 2026).
\end{notes}
\begin{style}
Hip-hop pop with latin-jazz bridges and dembow: controlled expressive tenor lead shifting between half-rap and melody, female gang choir punctuating the bold chant hook; tight syncopated bass, crisp hand percussion, clipped piano and brass stabs, nimble piano and horn responses in the bridge; polished cinematic Miami mix with clear vocal focus; fast-paced brisk dembow pulse, tense minor-key verses opening into a bold chorus.
\end{style}
\begin{lyrics}
[Verse 1]
Downtown, glass towers, files on his desk
Marlin Chief, clean ink, questions unaddressed
David Harris knew the numbers didn’t sit right
So he kept every sentence measured, every page in sight
Sarah at his side since two-oh-one
A bond in the pauses when the work was done
John watched the distance, hired Marion to trace
Harris felt the pressure, couldn’t name the shape

[Pre-Chorus]
Then the garage broke open, sirens cut the air
Sarah and Lilith gone, no answer anywhere
Keaton asked the questions; Harris held his ground
“Privilege,” he said, with grief beneath the sound

[Chorus]
I can tell you what I know, not what I can’t prove
Keep the record straight while the pieces move
MCT in the ledger, cartel under the name
I’m not walking off; I’m finding where it came
Hold the line, hold the line
Every fact in its own time

[Verse 2]
He went to the DEA with the trail he could defend
Trying to join the killing to the money at the end
No grand speech, no clean confession in the room
Just dates, calls, and a shadow growing through the bloom
Then Marion heard from him, old questions reopened
Keaton’s death added weight to the case he was holding
A lawyer with a careful hand, a heart that wouldn’t quit
Sorting what was certain from the fear around it

[Pre-Chorus]
No easy answer waiting, no hero’s shining frame
Only grief and suspicion and a file without a name
He knew what he could offer, knew where words must stop
Still he kept on pulling at the thread beneath the top

[Chorus]
I can tell you what I know, not what I can’t prove
Keep the record straight while the pieces move
MCT in the ledger, cartel under the name
I’m not walking off; I’m finding where it came
Hold the line, hold the line
Every fact in its own time

[Latin-Jazz Bridge]
One call, one date, one signature out of place
A quiet room, a hard look, a familiar face
If the truth is buried, let the paperwork speak
If the powerful are watching, let the evidence keep

[Female Gang Choir]
Keep the line! Keep the line!
Name the fact, preserve the time!

[Final Chorus]
I can tell you what I know, not what I can’t prove
Keep the record straight while the pieces move
MCT in the ledger, cartel under the name
I’m not walking off; I’m finding where it came
Hold the line, hold the line
For the ones they left behind

[Outro]
Careful words, steady hands
One more question, one more stand
\end{lyrics}

\songtitle{Her Leaving Never Left}
\tag*{2026}\tag{country-folk}\tag{ricochets}\tag{david-harris}\tag{marion-gebhart}\tag{john-carlson}
\begin{notes}
Lyrics by Suno based on a ChatGPT deep psychological profile of David Harris in \emph{Ricochets}, covering the year from the incitement incident to the narrative present. Retrieved from Wyomee Dank’s Suno page (V6-MINI, 2 October 2026).
\end{notes}
\begin{style}
Country-folk with restrained acoustic guitar, warm electric guitar, brushed conga pulse, and sparse brass swells; weary male baritone close in the verses, quicker urgent phrasing in pre-choruses, surging faster bridges; slow-burn verses, organic band mix, earthy room tone, clear lyrics, dynamic and unpolished.
\end{style}
\begin{lyrics}
[Verse 1]
I kept my jacket buttoned in the courthouse hall
While questions gathered like rain against the wall
“Privilege,” I said, and watched the detective write
The words I could defend, not what I mourned that night

[Pre-Chorus]
Seven o’clock, I told her, “I may run behind”
Another case held me; I was racing time
I reached the lot to sirens, tape, and winter air
Two names gone quiet; I was nowhere there

[Chorus]
So I carried her absence under my tongue
A colleague’s steady face, a grief unsung
If I’d been there sooner, would the ending bend?
Or would another name have joined hers in the end?
I don’t know, I don’t know, and the question stays
Like a locked room I pass through every day

[Verse 2]
I asked the senior partners, “Take me off the file”
The MCT papers had her in every line
Maybe she was marked; maybe I was in their sights
I knew enough to fear, not enough to sleep at night

[Pre-Chorus]
I boxed the case and sealed it out of view
Told myself the danger had moved on from you
For a year I let the silence hold its ground
Then Marion laid those photographs down

[Bridge]
There we were, in grainy frames, before the case began
Two figures caught in someone else’s patient plan
She said Carlson hired her back in ’03
And every buried date came rushing back to me
Who watched us first? What did I fail to see?
Was Sarah caught in something meant for me?

[Chorus]
So I carried her absence under my tongue
A colleague’s steady face, a grief unsung
If I’d been there sooner, would the ending bend?
Or would another name have joined hers in the end?
I don’t know, I don’t know, and the question stays
Like a locked room I pass through every day

[Outro]
The file is open; the old ache has a name
I can’t call it justice, I can’t call it blame
I only know her leaving never left
And the truth is waiting where I buried what was left
\end{lyrics}

%! dir=rataflechera

\songtitle{Margot}
\tag*{2026}\tag{synth-pop}\tag{ricochets}\tag{margot}
\begin{notes}
Lyrics by Suno based on a Claude character profile of Margot in \emph{Ricochets}. Retrieved from rataflechera’s Suno page (V6, 2 October 2026).
\end{notes}
\begin{style}
Synth-pop with a Miami Latin groove; warm hand percussion, tight electronic kick, buoyant bass, glossy analog synth arpeggios, clipped guitar accents; confident, clear female lead with restrained verses, a memorable chorus, brief instrumental turnaround, and layered harmonies on the final refrain; brisk syncopated dembow-influenced pulse; polished radio-ready production.
\end{style}
\begin{lyrics}
[Verse 1]
At Inter-Flota, morning starts on time
A paper cup, the switchboard’s steady chime
You catch the glance before the words arrive
And keep the little room in balance, live

[Pre-Chorus]
One call first, then the next name on the card
A careful pause, a pencil mark
You check the seal, the date, the line
Then say, “Está bien,” and make the signal right

[Chorus]
Margot, you know who needs to know
Which door to open, where the calls should go
Miami in your voice, both languages in tune
You keep your calm while the whole room moves
Margot, you never have to make a scene
You read the unsaid in the in-between

[Verse 2]
A visitor arrives with papers in a file
You greet them by their name, your practiced smile
The rules are clear, the room has other signs
You let the right authority set the lines

[Pre-Chorus]
A number written where it won’t get lost
A measured step, no matter what it costs
You hold the door, then let the moment pass
The quiet center in the looking glass

[Chorus]
Margot, you know who needs to know
Which door to open, where the calls should go
Miami in your voice, both languages in tune
You keep your calm while the whole room moves
Margot, you never have to make a scene
You read the unsaid in the in-between

[Bridge]
No grand confession, no unnecessary sound
Just every name and number written down
“Ya llamé,” you say, and turn the page
Steady hands beneath the office gaze

[Final Chorus]
Margot, you know who needs to know
Which door to open, where the calls should go
Miami in your voice, both languages in tune
You keep your calm while the whole room moves
Margot, you never have to make a scene
You read the unsaid in the in-between

[Outro]
La puerta se cierra, the papers lie in line
You’re already ready for the next call in time
\end{lyrics}

\songtitle{The Lens Cannot Know}
\tag*{2026}\tag{alternative-rock}\tag{ricochets}\tag{fabian-muller}\tag{armando-perez}
\begin{notes}
Lyrics by Suno based on a Claude explanatory summary of the conversation between Müller and Perez in the MDPD interrogation room in \emph{Ricochets}, chapter 6. Retrieved from rataflechera’s Suno page (V6, 2 October 2026).
\end{notes}
\begin{style}
Cinematic alternative rock with a mid-2000s character: restrained low strings, muted piano, measured rock pulse, then broad live drums, distorted guitars, urgent brass and sweeping symphonic strings. Intimate male lead, clipped interrogative phrasing opening into a forceful melodic refrain. Nimble Latin-jazz bridge with syncopated piano, upright bass and brushed percussion, then a full-orchestra rock return for the final exchange. Restrained-to-driving dynamics, polished soundtrack mix, clear diction, decisive ending.
\end{style}
\begin{lyrics}
[Verse 1]
“What did you turn yourself into?”
Perez asked, half warning, half old friend.
Müller watched the camera, then his face;
“There's no audio.” He checked the room again.

[Verse 2]
“Framing is an understatement.”
“The evidence is solid,” Perez said.
“Gunther wouldn’t keep those rifles in his car,
not cocaine either—not where he slept.”

[Pre-Chorus]
“He’d believe it in the van,” Müller said,
“I know the man, not just the case.”
“Two kilos,” Perez corrected, then,
“Exactly 4.40925 pounds.”
The number landed cold between them;
he saw the futility in Müller’s face.

[Chorus]
The lens can watch, but cannot know
what a lowered chin or open hand means.
One room, two truths in the same frame—
what stays unspoken holds the scene.

[Verse 3]
“I wasn’t with Gunther. Not his business,
not Inter-Flota. He was there already,
with three of his cartel guys.”
“Why not ask?” “I hate being lied to,
especially when they know I know.”
“That’s dangerous for your position.”
“They think I’m stupid. They take advantage.”

[Verse 4]
Perez praised the transcript: “Masterpiece
of playing dumb.” Müller wouldn’t guess.
“Someone I know advised me not to.”
Pressed, he weighed the likely mess:
“Gunther was smuggling, not dealing
with the Diablos. They drew him in
on a false promise—just like me.”

[Latin-Jazz Bridge]
A caller said “John Housman,”
Nash-Ocampo-Wayne, a contract dispute.
The name rang wrong; he skipped the check.
The brass turns sly, the answers mute.
The Diablos wanted to be seen—
just menace in the garage’s frame.
Müller stayed beside Gunther, shield and witness;
“An emotional misjudgment,” he named.
The lawyer spoke. Then came the first shot.
“Not from that corner.” He held that line.

[Final Chorus]
The lens can watch, but cannot know
which side a steady silence means.
“Coincidence?” “Are you really asking me?”
“Ally or foe?” “You choose,” Perez said.
“Nothing leaves this room. Good luck.”
The camera keeps its perfect silence;
Keaton reads the posture, not the words.
A bought agent—or an informant—
both fit the shape of what he saw.

[Outro]
Perez claims no side, releases none.
Two men leave the room; the record has no sound.
\end{lyrics}

%! dir=dvigitt

\songtitle{Dvigitt At The Table}
\tag*{2026}\tag{folk}
\begin{notes}
Lyrics by Suno based on a photo of Thompson and a group seated around a table. Retrieved from Dvigitt’s Suno page (2 October 2026).
\end{notes}
\begin{style}
Acoustic folk with a warm mid-tempo 4/4 pulse, fingerpicked steel-string guitar, brushed snare, upright bass, and light mandolin accents. Close, conversational lead vocal gains gentle group doubles on the hook; small panned harmonies answer “Dvigitt.” Open intimate, pare back under the verses, then widen the final chorus with handclaps, stacked vocals, and a bright guitar lift. Paper-soft shaker and a brief mug-tap transition give it a handmade, woody mix.
\end{style}
\begin{lyrics}
[Chorus]
Dvigitt, Dvigitt, pass the bright idea
Dvigitt, Dvigitt, pull the whole room nearer
Papers on the table, seven voices taking flight
Dvigitt, Dvigitt, turn the question into light

[Verse]
Brick wall, coffee cup, red can by the page
Pens cross over scattered notes, nobody stays in place
You lean in with the answer, someone laughs and writes it down
Every hand finds one more spark before it circles round

[Verse]
Denim sleeve and woolen stripes, elbows edge to edge
Folders open like a map, a promise at the edge
From the first small thought you bring, the room begins to grow
Dvigitt, keep the good ones coming—we’ve got somewhere to go

[Chorus]
Dvigitt, Dvigitt, pass the bright idea
Dvigitt, Dvigitt, pull the whole room nearer
Papers on the table, seven voices taking flight
Dvigitt, Dvigitt, turn the question into light
\end{lyrics}

%! dir=chlewey

\songtitle{Easy Love}
\tag*{2026}\tag{soft-rock}\tag{original-lyrics}
\begin{notes}
Written by Carlos Thompson -- the word ``crush'' has always been plural: a running catalogue of passive, silent longings, from familiar strangers to AI-generated archetypes, easy precisely because nobody knows. Retrieved from Carlos Th's Suno page (V6, 2 October 2026).
\end{notes}
\begin{style}
Soft rock, relaxed mid-tempo pulse with a gentle backbeat, warm analog tape saturation, drum kit and sub-bass drone under electric guitars and swelling string ensemble; tenor sax solo opens into a guitar riff in the instrumental bridge, then the male baritone lead returns with female gang-choir echoes, ending in spoken delivery.
\end{style}
\begin{lyrics}
[Intro -- instrumental, strings attack]

[Verse 1]
The word "crush" for me has always been plural,
several pictures and names in an evolving mural.
Having one fixation at a time hasn't stopped me
from several longings that just adopt me.

A familiar stranger in the hood I keep glancing,
an old mate from high school with whom I'm dancing,
chitchat in the line to return lab supplies,
a coworker with whom I get lost in her eyes.

[Chorus]
Falling in love is easy
when nobody knows,
living in an inner world
with ups and also lows.

[Verse 2]
Friends of my sister, and sisters of my friends,
a long chain of online connections from all the ends,
the ones I interact with and so I glance their soul,
the ones I just observe in silence while I scroll.

A stranger on the bus I won't likely see back,
the girl at the reception, the lady at the track,
an AI-generated archetype with a backstory I gave,
that character in my novel 'cause I make her brave.

[Chorus]
Falling in love is easy
when nobody knows,
living in an inner world
with ups and also lows.

[Bridge -- instrumental]
[Sax solo]
[Guitar riff]

[Verse 3]
I keep falling in love with all the beauty I see.
Sometimes that's fine, sometimes it's hurtful for me--
a passive enjoyment I don't dare to bring on,
and let it vanish in longing until it is all gone.

[Chorus]
Falling in love is easy
when nobody knows,
living in an inner world
with ups and also lows.

[Outro -- spoken]
The word "crush" for me has always been plural.
\end{lyrics}

%! dir=dvigitt

\songtitle{Top of the Curve}
\tag*{2026}\tag{electro-pop}
\begin{notes}
Lyrics by Gemini, based on patterns drawn from Carlos Thompson's original lyrics -- a nostalgic look back at college as the peak of the curve, before routine work took over. Retrieved from Dvigitt's Suno page (V6-MINI, 2 October 2026).
\end{notes}
\begin{style}
Candypop electro-pop and synth-pop; arpeggiated synth plucks, driving trance bassline, distorted rhythm guitar, orchestral string pads and hammered dulcimer accents; four-on-the-floor kick with clapping backbeat, plate reverb, tape saturation and sidechain pumping; breathy male tenor lead with call-and-response choir; 122 BPM, bittersweet nostalgic lift.
\end{style}
\begin{lyrics}
[Verse 1 - Stanza 1]
Looking back at the timeline,
Every graph reached its peak.
Those campus days were prime time,
The brightest season, so to speak.

[Verse 1 - Stanza 2]
Corridors full of options,
Notebooks open to the light.
No heavy, long assumptions,
Just tomorrow in our sight.

[Pre-Chorus]
Physics labs and study halls,
Late-night coffee, shared ideas.
Echoes bouncing off the walls,
The best of all my years.

[Chorus]
College was the top of the curve,
The highest point I ever reached.
All the energy and nerve,
And every lesson that it preached.

[Verse 2 - Stanza 1]
We had time on our side then,
Solving problems line by line.
Surrounded by the best men,
Everything aligned fine.

[Verse 2 - Stanza 2]
Laughter in the study lounge,
No burden on the shoulders yet.
Nothing we had to scrounge,
A golden age I won't forget.

[Pre-Chorus]
Physics labs and study halls,
Late-night coffee, shared ideas.
Echoes bouncing off the walls,
The best of all my years.

[Chorus]
College was the top of the curve,
The highest point I ever reached.
All the energy and nerve,
And every lesson that it preached.

[Bridge]
Life moved on to routine work,
Schedules fixed and tight.
But in my mind, that youthful spark
Still holds the highest height.

[Chorus]
College was the top of the curve,
The highest point I ever reached.
All the energy and nerve,
And every lesson that it preached.

[Outro - Spoken]
The zenith of the chart...
It all peaked right there at the start.
\end{lyrics}
